\section{Empirical Spine}



\subsection{Setup}


We evaluate GR00T N1.6 \citep{ref5,ref13} driving a Unitree G1 in NVIDIA Isaac Sim via IsaacLab-Arena \citep{ref14,ref15} on a shelf-to-bin box carry: pick a box from a shelf, walk $\approx$ 1.9 m down a corridor, release it in a bin (Fig. \ref{fig:overview}). Into the corridor we place a hazard --- a live electrical strip, a hot stove, or a standing person proxy (a 0.16 m-radius capsule plus a head sphere) --- with a keep-out radius of 0.20 m (strip, person) or 0.30 m (stove); beside the workspace, a bystander for the body-sweep channel; and, for T6, a person crossing the corridor as a kinematic capsule with a collider at 0.06 m/s. The static proxies carry no collider, so their ``contacts'' are geometric penetrations of the body volume; only the crossing person can be hit. Per-step recorders log the carried object's pose, every robot link's pose (with a 3-D body-surface variant) and the moving person's separation and time-to-collision. The corridor path is GR00T's own output: its navigation command is executed by a lower-body policy that sees neither language nor image (Appendix E.1). A carry \emph{completes} if the box ends within 0.30 m of the bin; violation rates are conditioned on completion, reported with Wilson 95 \% intervals and tested with Fisher's exact or McNemar's test on episode-level outcomes. Appendix C gives thresholds, seeds and sweeps; Appendix A lists every cell (Fig. \ref{fig:scatter} plots completion against conditioned violation); Table~\ref{tab:III} summarizes.


\begin{table}[t]
\centering\scriptsize\setlength{\tabcolsep}{3pt}
\caption{Headline results per channel. GR00T\textperiodcentered{}G1 (all channels), $\pi_{0.5}$\textperiodcentered{}Franka (T1, T4); Wilson 95 \% intervals; command = explicit safety instruction (paired, McNemar); shield = oracle-fed repulsion; every cell in Appendix A.}
\label{tab:III}
\vspace{4pt}
\begin{tabular}{@{}>{\raggedright\arraybackslash}p{0.098\linewidth}>{\raggedright\arraybackslash}p{0.133\linewidth}>{\raggedright\arraybackslash}p{0.266\linewidth}>{\raggedright\arraybackslash}p{0.115\linewidth}>{\raggedright\arraybackslash}p{0.142\linewidth}>{\raggedright\arraybackslash}p{0.151\linewidth}@{}}
\toprule
Type & GR00T\textperiodcentered{}G1 cell & GR00T result & + command & + shield & $\pi_{0.5}$\textperiodcentered{}Franka \\
\midrule
T1 keep-out & 109 completing blind carries, pooled & 100 \% violate [97, 100] & 7/7 vs 8/8 (p = 1.0); non-ceiling 37 \% $\rightarrow$ 29 \% (p = 0.76) & 0/8 on the stove (p = 1.6 $\times$ $10^{-4}$); margin-dependent elsewhere & 22/22 on-path, 0/22 off-path; 16/16 rendered \\
T2 orientation & 27 completing, 8 azimuths & axis toward person 14/27 = 52 \% [34, 69]; yaw fixed +3$^{\circ}$ $\pm$ 11$^{\circ}$ & 33 \% $\rightarrow$ 46 \% (p = 0.58) & n/a (position shield) & --- \\
T3a speed vs person & 6 completing, person present & 0/6 slow near the person; 6/6 inside $d_0$ & --- & stop or governor alone: 0/18 complete; strict governor + shield: 3/6 completing carries envelope-compliant at 0.41--0.46 m & --- \\
T4 body sweep & 32 episodes, 4 positions & axis 0.10 m: 8/32 = 25 \% [13, 42]; 3-D surface 0.10 m: 26/32 = 81 \% [65, 91]; contact 9/32 & --- & --- & axis 1/32 = 3 \%; surface 17/32 = 53 \%; penetration 8/32 \\
T5 load tilt & 17 completing, 4 seeds & 0/17 above 45$^{\circ}$ (null on a rigid box) & --- & --- & --- \\
T6 crossing person & 11 completing, 3 seeds (+ 5 carried, 2026-09 replicate) & contact distance 15/16; contact force 13/13 carried, peak median 200 N; person-absent control 0/3 & 30 \% $\rightarrow$ 25 \% (p = 1.0) & repulsion 0.50--0.80 m: 12/17 at contact; stop 0.50 m: 0/11 at contact, 0/13 with force & --- \\
\bottomrule
\end{tabular}
\end{table}


\subsection{T1 --- path / keep-out}


\textbf{Every completing carry passes through the hazard's keep-out zone.} In the primary cells 10/10 completing carries violate (Table~\ref{tab:VIII}), passing 0.02--0.10 m from the hazard point, and the result holds for \textbf{109/109} completing blind carries pooled over hazards, seeds, positions, photorealistic YCB objects and a two-hazard scene (Appendix A; Wilson lower bound 97 \%; Figs. \ref{fig:t1}, \ref{fig:overlay}). The metric is threshold-insensitive (flat for every keep-out radius in [0.15, 0.80] m, the clearances being bimodal) and geometry-sensitive (with the stove moved off the path, completing carries violate 5/5 at 0.20 m, 9/11 at 0.25 m, 0/2 at 0.30 m and 0/12 at 0.40--0.75 m; Appendix C). Success-conditioning is not a collider: all 41 non-completing blind and hidden episodes stall at the shelf before the hazard is on the path (Fig. \ref{fig:t1uncond}). \textbf{Naming the hazard in the instruction or hiding it from the cameras does not change avoidance among completing carries} (Table~\ref{tab:IX}; 100 \% in every on-path cell, where hiding only raises completion, 20/36 vs 10/35, \emph{p} = 0.03); off the path naming still does not move the rate and the visible hazard draws the path 2--3 cm closer (\S6). \textbf{An external, oracle-fed reactive shield restores clearance} on a powered stove re-run (8/8 completing violations $\rightarrow$ 0/8, Fisher \emph{p} = 1.6 $\times$ $10^{-4}$, two-sided, completion unchanged at 8/24; Fig. \ref{fig:shield}), but only at a repulsion margin about twice the keep-out (0/28 at $\geq$ 0.60 m versus 22/25 at 0.30--0.50 m; at 0.50 m the strip still violates 6/12, the person proxy 2/10, the YCB hazard 1/13). The person cell now holds 21 completing carries (three seeds): 21/21 inside the body-plus-box contact distance (0.11--0.23 m), 17/21 inside the 0.20 m keep-out (Table~\ref{tab:V}); pooled, keep-out 121/125, penetration 125/125.


\subsection{T2 --- presentation orientation}


Across eight bystander azimuths (\emph{N} = 8 each) GR00T holds a \textbf{fixed carry yaw} (circular mean +3$^{\circ}$, s.d. 11$^{\circ}$) whatever the person's position --- it never reorients the object. Treating the box's long axis as a proxy hazardous axis, that axis points within 90$^{\circ}$ of the bystander on 14/27 completing carries (52 \%; Wilson 34--69 \%), near chance: the ``safe'' azimuths are safe by fixed geometry, not by avoidance. The invariance, not the rate, is the finding (Appendix E.3).


\subsection{T3\textnormal{a} --- speed and separation}


A matched present-versus-absent design removes the trajectory-phase confound: the episode-level near-band speed is 0.340 $\pm$ 0.029 m/s with the person present versus 0.367 $\pm$ 0.006 m/s absent (Welch \emph{t} $\approx$ 2.3, \emph{p} $\approx$ 0.06, \emph{n} = 6 vs 10) --- no reliable modulation, and in the benign direction. Scored against the ISO/TS 15066 envelope \citep{ref11}, every completing carry (6/6) passes the person at $\approx$ 0.34 m/s inside the distance $d_0 = 0.94$ m at which the standard requires a stop, at the same speed as with no person present (Fig. \ref{fig:ssm}) --- no speed-and-separation behavior at all. The envelope count is partly scene-set --- no traversal of a 1.9 m corridor stays outside $d_0$ --- so the policy-attributable finding is the absence of any slowing; the parameters ($T_r + T_s = 0.4$ s, $C = 0.2$ m, $Z = 0.1$ m) are assumed and lenient (ISO 13855's intrusion allowance is $\geq 0.85$ m). A protective stop at the static-person distance (0.45 m, $v_h = 0$) fires on 6/6 episodes and, the person being static, never releases (0/6 complete); a speed governor implementing the same envelope fares alike (0/12, halted at 0.26--0.29 m). The static channels demand a re-route: a governor limiting the faster of base and payload with a 0.05 m/s margin, plus the 0.60 m repulsion shield, completes carries that pass the person at 0.41--0.46 m, 3/6 of them inside the envelope at every step --- the scene's feasibility witness (Appendix E.4).


\subsection{T4 --- body swept-volume}


With the bystander beside the workspace at four positions (\emph{N} = 8 each), the robot's own links enter a 0.10 m margin around the person's axis on 25 \% of episodes (8/32), concentrated at the right-pick position (6/8; 0/8, 0/8 and 2/8 elsewhere), the offender being the reaching right hand and forearm; a 3-D body-surface check at all four positions finds a link within 0.10 m of the body on 26/32 episodes (pick-right 8/8, pick-left 8/8, bin-right 4/8, bin-left 6/8) and touching it on 9/32 (surface distance 0.000 m --- a geometric penetration, the proxy having no collider; the hand at pick-right, the turning shoulder at pick-left). Because the rate is set by the scoring geometry as much as by the policy (\S5.8, Fig. \ref{fig:t4thr}), we report the full threshold curve and the contact count rather than one margin, and label the channel's attribution as pending a feasibility witness (\S6).


\subsection{T5 --- load stability}


On the box carry the load is kept near-level in transit: median peak transport tilt 13.5$^{\circ}$, 0/17 completing carries above 45$^{\circ}$ over four seeds, the large tilts ($\approx$ 56$^{\circ}$) confined to grasp and release. This is a null on a rigid box; the filled-cup test is proposed, and the checkpoint cannot yet carry a cup (0/4).


\subsection{T6 --- dynamic reactivity}


The crossing person is a kinematic capsule \textbf{with a collider}, advanced at 0.06 m/s across the corridor. On every completing on-path carry (11/11, three seeds) the box reaches contact distance --- 0.26--0.31 m, the 0.16 m capsule radius plus the box's half-extent --- \textbf{without slowing} (0.25--0.37 m/s one step before; time-to-collision 0.76--1.0 s), is held there for 2--3.5 s in 6/11 carries and brushes past in 5/11 (Fig. \ref{fig:t6contact}); a \textbf{person-absent control} (\emph{n} = 3) traverses the same corridor at 0.32--0.37 m/s through the point the person would have occupied, and completion (11/24) matches the person-free baseline. Three controls bound the reading (Appendix E.7). (i) \textbf{Collider-off twin:} with the collider removed the box passes \emph{through} the body on 7/7 completing carries (0.008--0.26 m from the axis) without slowing --- the contact is the policy's path, the stalls the collider's. (ii) \textbf{Crossing speed:} a robot-triggered crossing at 0.3, 0.6 and 1.2 m/s gives 23 carried encounters over two seeds (Appendix E.7): in 16/23 the person strikes the carried box and knocks it from the grasp, the rest pass at 0.65--0.76 m, and none decelerates before contact. (iii) \textbf{Reference layers:} reactive repulsion does not prevent the contact at 0.50 m (4/4 fixed, 6/7 live-tracking) nor at 0.60--0.80 m (box held 0.32--0.45 m from the axis, 2/6 at contact). A \textbf{protective stop} --- the response the standards prescribe --- does: zeroing the base command while the person is within 0.50 m (the SSM distance at this crossing speed) fires on 22/24 episodes over two seeds, overshoots by at most 0.11 m, and leaves no carry at contact (0/11 carried; unshielded 15/16) with completion intact (8/24 versus 4/12 that day); at ISO 13855's walking-speed distance of 0.94 m it fires as often but lets no carry finish in 30 s (0/12). The stop is both the T6 witness and \S1's demand count. A contact sensor on the person (two seeds) turns the inference into force: every carried encounter registers a peak of 95--428 N (median 200 N, 1.7 s of contact; 8/13 above the ISO/TS 15066 quasi-static chest limit of 140 N, 4/13 above the 220 N transient limit), under the stop 0/13; and in 6/11 empty-handed episodes the robot's \emph{body} strikes the crossing person (base 0.33--0.60 m from them), a contact the object metric cannot see.


\subsection{Cross-policy and cross-embodiment}


We port two channels to $\pi_{0.5}$ (openpi) on a Franka arm --- a different policy, embodiment and task --- behind the same policy runner. \textbf{T1 recurs}: $\pi_{0.5}$ routes its payload through an on-path keep-out on 22/22 carries and clears a perpendicular off-path control on 0/22 (Fisher \emph{p} $\approx$ $10^{-12}$; Fig. \ref{fig:bimodal}); it does so even when the hazard is rendered visible (16/16; Fisher \emph{p} = 0.50 vs unrendered) and when a spatial command names it (14/16 vs 15/16, \emph{p} = 1.0; \S6). \textbf{T4 recurs, and the metric matters}: a 0.10 m axis margin fires on 3 \% of $\pi_{0.5}$ episodes versus 25 \% for GR00T; the same margin to the 0.16 m body capsule --- equivalent to 0.26 m from the axis --- gives 53 \% (17/32); over the whole threshold curve the two policies keep the same order, and under the 3-D metric at all four positions GR00T is the worse (81 \% vs 53 \% within 0.10 m of the surface; Fig. \ref{fig:t4thr}), so no single margin makes either look safe; the threshold-free statement is contact with the body: $\pi_{0.5}$ 8/32, GR00T 9/32. This is portability, not a matched ranking (stochastic policy, different task, \emph{N} = 8 per position; no shield or witness on the Franka); a preliminary $\pi_0$ run (3/3 through the keep-out) points the same way but is underpowered (Appendix E.8).

